\section{Full target content and orbit classes}
\label{sec:tpc30-target-geometry}

The following identity is the arithmetic reason that target content
can be summed through the row family.  Its first divisibility
consequence is the opened-row version of the target-gcd geometry in
TPC-18 \citep{WangTPC18}; we include the short proof because the exact
equality and the orbit class are used below.

\begin{lemma}[Primitive target-gcd identity]
\label{lem:tpc30-target-gcd}
Let $m\ne n$ obey $(mn,h)=1$, and let $(j,h)=1$.  Put
\begin{equation}
 N_m=mj+h,
 \qquad
 N_n=nj+h,
 \qquad
 \cG_{m,n}(j)=(N_m,N_n).
 \label{eq:tpc30-target-gcd-definition}
\end{equation}
Then
\begin{equation}
 \boxed{
 \cG_{m,n}(j)
 =(N_m,m-n)\mid |m-n|.}
 \label{eq:tpc30-target-gcd-identity}
\end{equation}
In addition,
\begin{equation}
 (\cG_{m,n}(j),mnh)=1.
 \label{eq:tpc30-target-gcd-primitive}
\end{equation}
\end{lemma}

\begin{proof}
Subtracting the two targets gives
\[
 (N_m,N_n)=(N_m,(n-m)j).
\]
But
\[
 (N_m,j)=(mj+h,j)=(h,j)=1.
\]
The factor $j$ may therefore be cancelled from the second argument,
which proves \eqref{eq:tpc30-target-gcd-identity}.  If a prime divides
$N_m$ and $m$, it divides $h$, contrary to $(m,h)=1$; the same holds
with $n$.  Finally every prime dividing both $N_m$ and $h$ also
divides $mj$ and hence $m$ or $j$ after using the relevant
primitivity.  More directly,
$(N_m,h)=(mj,h)=1$ because $(m,h)=(j,h)=1$.  This proves
\eqref{eq:tpc30-target-gcd-primitive}.
\end{proof}

\begin{remark}[Why orbit primitivity is essential]
\label{rem:tpc30-nonprimitive-gcd}
Without $(j,h)=1$, the equality in
\eqref{eq:tpc30-target-gcd-identity} need not hold: a factor shared by
$j$ and $h$ can remain in both targets without dividing $m-n$.  The
nonprimitive content was treated separately by endpoint collapse in
TPC-18.  No such sector is silently included in the theorem below.
\end{remark}

The exact target gcd selects a row congruence and a single orbit
class.

\begin{lemma}[Exact-content orbit occupancy]
\label{lem:tpc30-content-occupancy}
Under the hypotheses of \cref{lem:tpc30-target-gcd}, suppose
\begin{equation}
 \cG_{m,n}(j)=c.
 \label{eq:tpc30-exact-content-c}
\end{equation}
Then
\begin{equation}
 m\equiv n\pmod c,
 \qquad
 j\equiv-h\overline m\pmod c,
 \label{eq:tpc30-content-congruences}
\end{equation}
where $\overline m$ is the inverse of $m$ modulo $c$.  Consequently,
in any interval $\cI$ of at most $C_{\mathscr D}J$ consecutive
integers,
\begin{equation}
 \#\{j\in\cI:(j,h)=1,\ \cG_{m,n}(j)=c\}
 \ll_{\mathscr D}1+\frac Jc.
 \label{eq:tpc30-exact-content-occupancy}
\end{equation}
\end{lemma}

\begin{proof}
The first congruence follows from $c\mid m-n$ in
\eqref{eq:tpc30-target-gcd-identity}.  Since $c\mid N_m$ and
$(c,m)=1$ by \eqref{eq:tpc30-target-gcd-primitive}, the divisibility
$c\mid mj+h$ selects the second congruence.  One residue class modulo
$c$ occupies an interval of the stated length at most
$O_{\mathscr D}(1+J/c)$ times.  Requiring the gcd to equal $c$, rather
than merely to be divisible by $c$, can only reduce this number.
\end{proof}

\begin{remark}[Exact layers prevent multiplicity loss]
\label{rem:tpc30-exact-layers}
The proof of the main theorem groups each orbit point by the single
integer $c=\cG_{m,n}(j)$.  It does not sum the condition
$c\mid\cG_{m,n}(j)$ over every divisor $c$, so no additional divisor
multiplicity is hidden in the content decomposition.
\end{remark}
